\documentclass[conference]{IEEEtran}
\usepackage{amsmath,amssymb,newtxmath}
\usepackage[no-math]{fontspec}
\setmainfont{Times New Roman}
\usepackage{xeCJK}
\setCJKmainfont[AutoFakeBold=2]{SimSun}
\usepackage{graphicx,booktabs,array,tabularx,seqsplit,xurl,hyperref}
\hypersetup{hidelinks,unicode=true,pdftitle={Research manuscript template},pdfauthor={Author name}}
\setlength{\emergencystretch}{2em}
\raggedbottom
\renewcommand{\arraystretch}{1.15}
\begin{document}
\title{Research manuscript template}
\author{\IEEEauthorblockN{Author name}\IEEEauthorblockA{Institution}}
\maketitle
\begin{abstract}
Summarize the question, method, principal findings and scope.
\end{abstract}
\begin{IEEEkeywords}
modeling; validation
\end{IEEEkeywords}
\section{Methods}
\hypertarget{chapter-01}{}
Use SI units and RMS phasors. Reference [1] defines the special functions used by the round-wire solution.

\begin{equation}\delta=\frac{1}{\sqrt{\pi f\mu\sigma}}\label{eq:skin}\end{equation}
\begin{table*}[!htbp]
\centering
\caption{Symbols and units}
\begingroup\footnotesize
\begin{tabularx}{\textwidth}{>{\raggedright\arraybackslash}X>{\raggedright\arraybackslash}X}
\toprule
Symbol & Unit\\
\midrule
f & Hz\\
σ & S/m\\
\bottomrule
\end{tabularx}\endgroup
\end{table*}
\section{Results and discussion}
\hypertarget{chapter-02}{}
Replace the illustration with traceable results and explain comparison conditions, errors and conclusions.

\renewcommand{\refname}{References}
\begin{figure}[!htbp]\centering\fbox{\rule{0pt}{22mm}\rule{.8\linewidth}{0pt}}\caption{Figure placement example; replace with actual results.}\label{fig:example}\end{figure}
\begin{thebibliography}{99}
\bibitem[1]{BESSEL} NIST. NIST Digital Library of Mathematical Functions, Modified Bessel Functions. 2026. \url{https://dlmf.nist.gov/10.25}
\end{thebibliography}
\end{document}
